\section{总结与延伸}

\subsection{讲者的核心总结}

\begin{figure}[H]
\centering
\includegraphics[width=0.95\textwidth]{slides-images/slide-050.jpg}
\caption{全课总结：如何在实际中进行 Scaling\protect\footnotemark}
\end{figure}
\footnotetext{Slides 第51页。}

Tatsu Hashimoto 在课程结尾总结了 Scaling Laws 实践的三个层面：

\begin{enumerate}
    \item \textbf{超参数设置}：通过 scaling laws 或 MuP 来确定学习率和 batch size，避免在大规模训练上进行代价高昂的网格搜索
    \item \textbf{稳定性策略}：使用 MuP 或 assume stability（假设超参数不随规模变化）来简化 scaling 流程
    \item \textbf{高效数据分析}：使用 WSD 学习率调度来降低 Chinchilla 式 data scaling 分析的计算成本
\end{enumerate}

\subsection{全课知识图谱}

本课建立了从理论到实践的完整认知链：

\begin{center}
\begin{tikzpicture}[node distance=0.8cm, every node/.style={draw, rounded corners, minimum width=3.5cm, minimum height=0.7cm, font=\small}]
\node (theory) {Scaling Law 理论基础};
\node (case) [below=of theory] {案例研究};
\node (mup) [below=of case] {MuP 理论与实证};
\node (practice) [below=of mup] {实践 Scaling 流程};
\draw[->, thick] (theory) -- (case);
\draw[->, thick] (case) -- (mup);
\draw[->, thick] (mup) -- (practice);
\node[draw=none, right=1cm of theory, text width=5.5cm, font=\scriptsize] {Chinchilla、isoFLOP、Critical Batch Size};
\node[draw=none, right=1cm of case, text width=5.5cm, font=\scriptsize] {Cerebras GPT、MiniCPM、DeepSeek、Llama 3};
\node[draw=none, right=1cm of mup, text width=5.5cm, font=\scriptsize] {激活稳定性、更新稳定性、$1/\text{fan\_in}$ 学习率};
\node[draw=none, right=1cm of practice, text width=5.5cm, font=\scriptsize] {小模型搜索 + MuP 迁移 + WSD 调度};
\end{tikzpicture}
\end{center}

\subsection{关键 Takeaways}

\begin{importantbox}{七条核心原则}
\begin{enumerate}
    \item \textbf{Chinchilla 分析是最可靠的 Scaling 工具}：isoFLOP 分析在所有团队、所有复现中都表现出色，是 scaling 的``黄金标准''
    \item \textbf{20:1 只是起点}：实际最优 token/parameter 比例因架构、数据质量、优化策略而异，可以显著高于 20:1
    \item \textbf{超参数 Scaling 总是嘈杂的}：学习率、batch size 的 scaling law 不如 loss scaling law 干净，但获取正确的数量级即可
    \item \textbf{WSD 是实用的学习率调度}：它使 Chinchilla 式分析几乎可以在一次训练中完成
    \item \textbf{MuP 解决了学习率迁移问题}：通过合理的初始化和每层学习率缩放，可以在小模型上搜索最优学习率并直接迁移
    \item \textbf{MuP 不是万能的}：它对 weight decay、非标准优化器、可学习 gain 敏感
    \item \textbf{认真的 Scaling 分析是区分成功与失败的关键}：DeepSeek 等模型的成功与其仔细的 scaling 研究密不可分
\end{enumerate}
\end{importantbox}

\subsection{拓展阅读}

\begin{itemize}
    \item Cerebras GPT 论文：\url{https://arxiv.org/abs/2304.03208}
    \item MiniCPM 技术报告：\url{https://arxiv.org/abs/2404.06395}
    \item DeepSeek LLM 论文：\url{https://arxiv.org/abs/2401.02954}
    \item Llama 3 论文：\url{https://arxiv.org/abs/2407.21783}
    \item Yang et al., Tensor Programs V (MuP)：\url{https://arxiv.org/abs/2203.03466}
    \item 《A Large-Scale Exploration of $\mu$-Transfer》：系统验证 MuP 鲁棒性的 ablation 研究
    \item 《A Practitioner's Guide to MuP》：MuP 的实践入门指南
    \item Hoffmann et al., Training Compute-Optimal Language Models (Chinchilla)：\url{https://arxiv.org/abs/2203.15556}
    \item 过度训练惩罚分析（UW + Apple）：估算超越 Chinchilla 比例时的 loss 退化
\end{itemize}

\end{document}
